\section{Jordan algebras and JB-algebras}\label{sec:jordan}

This section fixes the algebraic vocabulary used throughout the report. Every
object here is standard; we follow the axioms in Hanche-Olsen and
St{\o}rmer~\cite{HancheOlsenStormer1984} so that later approximate versions
(\Cref{sec:epsjb}) can be read as quantitative relaxations of these exact
identities. We work over a field of characteristic not $2$; for the analytic
results the scalars are real, as indicated.

\subsection{Jordan algebras}

The Jordan product symmetrises an associative product: for operators on a Hilbert
space, $a\jp b=\tfrac12(ab+ba)$ is self-adjoint whenever $a$ and $b$ are, so it
keeps us inside $\Bsa$ where the associative product $ab$ does not. Abstracting
the two identities that this symmetrised product satisfies gives the notion of a
Jordan algebra.

% PROV: HancheOlsenStormer1984 line 812 (eq 2.17) and lines 963-967 (item 2.4.1, eqs 2.29-2.30) -- Jordan product and the two defining axioms.
\begin{definition}[Jordan algebra]\label{def:jordan-algebra}
\status{(cited)}
For elements $a,b$ of an associative algebra, the \emph{Jordan product} is
\begin{equation}\label{eq:jordan-product}
  a \jp b = \tfrac12(ab+ba).
\end{equation}
An algebra $A$ with a product written $(a,b)\mapsto a\jp b$ is called a
\emph{Jordan algebra} if the following two identities hold for all $a,b\in A$:
\begin{align}
  a \jp b &= b \jp a, \label{eq:jordan-comm}\\
  a \jp (b \jp a^2) &= (a \jp b) \jp a^2. \label{eq:jordan-identity}
\end{align}
Identity~\eqref{eq:jordan-comm} states that $A$ is commutative;
\eqref{eq:jordan-identity} is the \emph{Jordan identity}. Here $a^2 := a\jp a$.
\end{definition}

% PROV: HancheOlsenStormer1984 line 812 and lines 965-967; notation harmonised to \jp.

\begin{remark}[Form of the Jordan identity]\label{rem:jordan-identity-form}
\status{(cited)}
The Jordan identity~\eqref{eq:jordan-identity}, written
$a\jp(b\jp a^2)=(a\jp b)\jp a^2$ in~\cite{HancheOlsenStormer1984}, is the same as
the symmetric-looking form
\[
  \bigl((a\jp a)\jp b\bigr)\jp a \;=\; (a\jp a)\jp(b\jp a),
\]
obtained by writing $a^2=a\jp a$ and applying
commutativity~\eqref{eq:jordan-comm} to each product. We use whichever form is
more convenient; they are equivalent in any commutative algebra.
\end{remark}

A Jordan algebra need not be associative. Its one indispensable substitute for
associativity is that any single element, together with its powers, generates an
associative subalgebra. This is what makes a functional calculus possible later.

% PROV: HancheOlsenStormer1984 lines 1013-1025 (items 2.4.4-2.4.5) -- power-associativity: powers defined inductively and the subalgebra on one element is associative.
\begin{proposition}[Power-associativity]\label{prop:power-associative}
\status{(cited)}
Let $A$ be a Jordan algebra and $a\in A$. Define powers inductively by $a^1=a$
and $a^{n+1}=a\jp a^n$ (and $a^0=\Id$ if $A$ is unital). Then the subalgebra
generated by $a$ is associative; equivalently, for all natural numbers $m,n$,
\begin{equation}\label{eq:power-assoc}
  a^{m+n} = a^m \jp a^n .
\end{equation}
Consequently every power $a^n$ is unambiguously defined.
\end{proposition}

% PROV: HancheOlsenStormer1984 lines 1013-1017; notation harmonised to \jp.

\subsection{JB-algebras}

To do analysis one adds a complete norm compatible with the product. The two
extra norm conditions below are the Jordan analogue of the $C^*$-identity: the
first fixes the norm of a square, the second forces sums of squares to dominate
their summands and is what makes the order structure (positivity = being a
square) well-behaved.

% PROV: HancheOlsenStormer1984 lines 2308-2316 (items 3.1.3-3.1.4) -- Jordan Banach algebra and the two JB-algebra norm conditions.
\begin{definition}[JB-algebra]\label{def:jb-algebra}
\status{(cited)}
A \emph{Jordan Banach algebra} is a real Jordan algebra $A$ equipped with a
complete norm satisfying
\begin{equation}\label{eq:submultiplicative}
  \norm{a \jp b} \le \norm{a}\,\norm{b}, \qquad a,b\in A.
\end{equation}
A \emph{JB-algebra} is a Jordan Banach algebra $A$ whose norm satisfies, in
addition, the following two conditions for all $a,b\in A$:
\begin{align}
  \norm{a^2} &= \norm{a}^2, \label{eq:jb-square}\\
  \norm{a^2} &\le \norm{a^2 + b^2}. \label{eq:jb-sum-of-squares}
\end{align}
If $A$ is unital with unit $\Id$, then~\eqref{eq:jb-square} gives $\norm{\Id}=1$.
\end{definition}

% PROV: HancheOlsenStormer1984 lines 2308-2316; source drops the hyphen in "JB-algebra" by a typesetting error noted at line 13; we restore it.

The motivating example is concrete: for a complex Hilbert space $\HH$, the
self-adjoint operators $\Bsa$ form a JB-algebra under~\eqref{eq:jordan-product}
and the operator norm. A norm-closed Jordan subalgebra of some $\Bsa$ is called a
\emph{JC-algebra}~\cite[3.1.2]{HancheOlsenStormer1984}, and every JC-algebra is a
JB-algebra. The content of the finite-dimensional theory is that, with one
exception, the converse also holds.

\subsection{Finite-dimensional classification}

In finite dimensions the JB-algebras are exactly the classical objects of
Jordan, von Neumann and Wigner~\cite{JordanVonNeumannWigner1934}: a real Jordan
algebra $A$ is called \emph{formally real} (or Euclidean) if
$\sum_{i=1}^n a_i^2 = 0$ in $A$ forces $a_1=\dots=a_n=0$. A finite-dimensional
unital formally real Jordan algebra is a JB-algebra in its order
norm~\cite[3.1.7]{HancheOlsenStormer1984}, and conversely every JB-algebra is
formally real~\cite[3.3.8]{HancheOlsenStormer1984}; so in finite dimensions the
two classes coincide. Their simple pieces are completely classified.

% PROV: HancheOlsenStormer1984 line 2254 (Theorem 2.9.6) together with line 2264 (2.9.7) and Theorem 2.9.8 (line 2270) -- classification of finite-dimensional simple formally real Jordan algebras; attributed to Jordan-von Neumann-Wigner [69] at line 2293.
\begin{theorem}[Jordan--von Neumann--Wigner classification]\label{thm:jnw-classification}
\status{(cited)}
Every finite-dimensional, formally real, unital Jordan algebra is a direct sum of
simple ones. Each simple summand is isomorphic to exactly one of
\begin{itemize}
  \item the one-dimensional algebra $\R$;
  \item a \emph{spin factor} $\R\Id\oplus H$, where $H$ is a real inner-product
    space of dimension $\ge 2$ (the case of two minimal projections summing
    to~$\Id$);
  \item a Hermitian matrix algebra $H_n(\R)$, $H_n(\C)$, or $H_n(\mathbb{H})$
    with $n\ge 3$, where $\mathbb{H}$ denotes the quaternions; or
  \item the \emph{Albert algebra} $H_3(\mathbb{O})$ of $3\times 3$ Hermitian
    matrices over the octonions $\mathbb{O}$.
\end{itemize}
Of these, only $H_3(\mathbb{O})$ is \emph{exceptional}: it is not isomorphic to a
Jordan subalgebra of self-adjoint operators on any Hilbert space; all the others
are JC-algebras.
\end{theorem}

\smallskip\noindent\textbf{Agent B review correction.}
The one-dimensional summand $\R$ is included explicitly; it is the $n=1$ case in
Hanche-Olsen--St{\o}rmer Theorem 2.9.6 and is not covered by the listed spin
factors.

% PROV: HancheOlsenStormer1984 line 2254 (2.9.6), line 2264 (2.9.7), line 1856 (Corollary 2.8.5), and line 2306; attribution at line 2293; original source JordanVonNeumannWigner1934.

\begin{remark}[Status of the exceptional algebra]\label{rem:exceptional}
\status{(cited)}
The exceptionality of $H_3(\mathbb{O})$ traces back to the non-associativity of
the octonions~\cite[2.8.5]{HancheOlsenStormer1984}. It is the only obstruction to
representing a finite-dimensional formally real Jordan algebra by self-adjoint
operators, and it is the reason the structure theory of \Cref{sec:programme} must
allow for a Jordan, rather than associative, limit object.
\end{remark}
